\section{The exact finite cyclotomic action}
\label{sec:cubic-finite-action}

The explicit bound in AC-K18 can be strengthened to a finite
descent statement.  This identifies the exact arithmetic data
that would determine all the remaining weights, without
performing a class-group calculation in a degree-54 field.

\begin{theorem}[AC-K20: action of order nine and a finite matrix recipe]
\label{AC-K20}
Keep $M=X'$, $\delta=\omega-\omega^2$ and
$\varepsilon_3(\gamma)=4$ as in AC-K18.  Then
\[
                    \gamma^9=1\quad\hbox{on }M,
       \qquad \gamma^3\ne1\quad\hbox{on }M.          \tag{KF.1}
\]
For
\[
 E_n=\mathbb Q(\mu_{3^{n+1}},\sqrt[3]{269}),
 \qquad n\geq0,
\]
the logarithmic descent isomorphisms consequently give
\[
 M\simeq\widetilde{\operatorname{Cl}}_3(E_n)
                      \quad(n\geq2),                 \tag{KF.2}
\]
compatibly with the cubic action and the remaining cyclotomic
action.  Thus $E_2=\mathbb Q(\mu_{27},\sqrt[3]{269})$, of
degree $54$ over $\mathbb Q$, suffices.  It is the first layer
for which the quotient map
$M\to M/(\gamma^{3^n}-1)M$ is an isomorphism.

More explicitly, suppose that a certified presentation of its
degree-zero logarithmic class group and its cyclotomic action
is given by
\[
 C=\mathbb Z^s/D\mathbb Z^s,
 \qquad D=\operatorname{diag}(3^{r_1},\ldots,3^{r_s}),
 \qquad G:C\longrightarrow C.                         \tag{KF.3}
\]
Here $G$ must represent the field automorphism fixing
$\sqrt[3]{269}$ and sending $\zeta_{27}$ to
$\zeta_{27}^4$, with its induced action on logarithmic ideal
classes.  For $j=1,\ldots,9$, put
\[
 C_j=\operatorname{coker}_{\mathbb Z}
       \begin{bmatrix}D&4^jG-I_s\end{bmatrix}.
                                                               \tag{KF.4}
\]
Then the exact exponent sequence is recovered by
\[
       m_e=1+\log_3\#C_j
       \quad\text{whenever }e\equiv j\pmod9.         \tag{KF.5}
\]
Thus nine finite Smith normal forms of one certified module
with its action determine every weight.  The invariant factors
of $C$ without $G$ do not constitute the data in this recipe.
No such degree-54 arithmetic certificate is asserted here.
\end{theorem}
\begin{proof}
\proofstep{1. A bound on the kernel of multiplication by $\delta$.}
Use the generator $x$ and the conjugation components from
AC-K18:
\[
 M_+=\mathbb Z_3[[B]]x,\qquad
 U=M_-=\mathbb Z_3[[B]]\delta x,
 \qquad B=\gamma-1.
\]
Multiplication by $\delta$ maps $M_+$ onto $U$.
Let $H$ be its kernel.  Since $\delta^2=-3$, $3H=0$.
The endomorphism $B$ is nilpotent, as proved in AC-K18.
Write $d=\dim_{\mathbb F_3}(M_+/3M_+)\in\{6,7\}$.
Because $3M_+=-\delta U$, the length of $3M_+$ is at
most the length of $U$.  The exact sequence defining $H$
therefore gives
\[
 \begin{aligned}
 \dim_{\mathbb F_3}H
 &=\operatorname{length}_{\mathbb Z_3}M_+
        -\operatorname{length}_{\mathbb Z_3}U
 \\
 &=d+\operatorname{length}_{\mathbb Z_3}(3M_+)
        -\operatorname{length}_{\mathbb Z_3}U
 \leq d\leq7.
 \end{aligned}
                                                               \tag{KF.6}
\]
In particular $B^7H=0$.

By AC-K18, the nonzero group $U$ has order at most nine
and $U/BU\simeq\mathbb Z/3$.
There are precisely the following possibilities for its
cyclic $\mathbb Z_3[[B]]$-module structure:
\[
 \begin{array}{ll}
 U=\mathbb Z/3^a,&a=1\text{ or }2,\quad
       B\text{ acts as }b\in3\mathbb Z_3;\\
 U=\mathbb F_3[B]/B^2.&
 \end{array}                                                   \tag{KF.7}
\]
In the first line, take $b=0$ if $a=1$; if $a=2$,
the quotient $U/BU$ forces $v_3(b)=1$.
In the elementary group of order nine, cyclicity and the same
quotient force the single nilpotent block in the second line.
Every case in (KF.7) is killed by
\[
                 F=(1+B)^3-1=B^3+3B^2+3B.           \tag{KF.8}
\]
Indeed this is immediate in characteristic three; for
$\mathbb Z/9$ it follows from $3\mid b$.
Consequently $FM_+\subset H$.
Since $3H=0$ and $F$ acts on $H$ as $B^3$,
\[
       (\gamma^9-1)M_+=F^3M_+=B^6FM_+.             \tag{KF.9}
\]
The first equality follows by expanding
$(1+F)^3-1=F^3+3F^2+3F$; the last two terms vanish on
$M_+$ because $FM_+\subset H$.

\proofstep{2. The two possible small components.}
Suppose first that $U=\mathbb Z/3^a$ as in the first
line of (KF.7).  The vector $\delta x$ is an abelian
generator of $U$ and has annihilator $(3^a,B-b)$ over
$\mathbb Z_3[[B]]$.  Hence $H$ is generated by
\[
                    t=(B-b)x,\qquad z=3^a x.
\]
Both lie in $H$.  Moreover
$Bz=bz+3^at=0$, because $3H=0$ and $3\mid b$.
Polynomial division, followed by $3t=0$ and $3\mid b$,
gives the exact identity
\[
             Fx=F(b)x+B^2t.                          \tag{KF.10}
\]
For $a=1$ our choice $b=0$ makes the first term zero.
For $a=2$, $F(b)$ is divisible by nine, so that the first
term is a multiple of $z$.
It follows that $B^6Fx=B^8t=0$, using (KF.6).
Since $x$ generates $M_+$ and all endomorphisms commute,
(KF.9) proves $\gamma^9=1$ on $M_+$ in this case.

In the remaining case $U=\mathbb F_3[B]/B^2$,
the annihilator of $\delta x$ is $(3,B^2)$.
Thus $H$ is generated by
\[
                         t=B^2x,\qquad z=3x.
\]
The relation $3H=0$ gives $B^2z=3t=0$.
From (KF.8) one obtains
\[
                         Fx=Bt+Bz.
\]
Therefore $B^6Fx=B^7t+B^7z=0$, again by (KF.6).
Equation (KF.9) proves the assertion on $M_+$.
On $U$, already $\gamma^3=1$, so this proves
$\gamma^9=1$ on all of $M$.

For the strict assertion in (KF.1), reduce $M_+$ modulo
three.  This is the cyclic nilpotent module
$\mathbb F_3[B]/B^d$ with $d\geq6$.
Here $\gamma^3-1=B^3$ is nonzero.  Thus the action of
$\gamma$ on $M$ has order exactly nine.

\proofstep{3. Finite descent and the matrix computation.}
The same logarithmic descent theorem used in (KL.9), applied
over $E_n$, identifies
\[
 \widetilde{\operatorname{Cl}}_3(E_n)
                 =M/(\gamma^{3^n}-1)M.              \tag{KF.11}
\]
Here the cyclotomic infinite extension and its maximal split
unramified module are unchanged; only the cyclotomic subgroup
has generator $\gamma^{3^n}$.
The exact primary-source input is the Schneider description
stated in the proof of Jaulent--Michel, Theorem~3, p.~7 of
the author's version, already checked in AC-K14:
\url{https://www.math.u-bordeaux.fr/~jjaulent/Articles/J44.pdf}.
It is an equality of compact coinvariants, not a conjectural
formula for a characteristic ideal.
Equation (KF.1) proves (KF.2) and the assertion about the
first stable quotient.  The field degree is
$3\varphi(27)=54$, since the cubic radical extension is
disjoint from the cyclotomic extension over
$\mathbb Q(\omega)$, as used in AC-K14.

Under (KF.2), the positive-weight wild kernel is
\[
              W_e(E)=C/(4^eG-1)C.                   \tag{KF.12}
\]
This is (KL.9), with the generator normalized by
$\varepsilon_3(\gamma)=4$.
If $C$ is given by (KF.3), adjoining the relations
$(4^eG-I_s)v$ to the columns of $D$ is exactly the
integer presentation (KF.4).
AC-K18 gives $\lambda^5C=0$ and therefore $27C=0$.
Since $4^9\equiv1\pmod{27}$, the quotient depends only
on $e$ modulo nine.  Finally AC-K12 and AC-G26 identify
$\log_3\#W_e=a_e=m_e-1$, which proves (KF.5).
Every step of the matrix recipe is finite exact algebra once
the stated arithmetic module and its action are certified.
\end{proof}

\paragraph{The remaining arithmetic certificate.}
A computation for this purpose must certify the degree-zero
logarithmic class module at three and the action induced by
$\zeta_{27}\mapsto\zeta_{27}^4$; an ordinary class group or
an unlabelled list of logarithmic invariant factors is
insufficient for (KF.4).  The induced action can be supplied
on a presentation by logarithmic ideal classes, with the images
of its generators verified modulo the certified relations.
The cubic action is a useful compatibility check, but its
matrix is not needed merely to take the orders in (KF.5).
Reversing the cyclotomic generator without changing the Tate
twist would permute the weight classes, so its normalization
is part of the certificate.  No uncontrolled degree-54
calculation is included in this reduction.

\paragraph{Finite-cover and even K-theory interpretation.}
The cubic and cyclotomic actions on $M$ factor through
$\operatorname{Gal}(E_2/K)\simeq C_3\times C_9$.
Thus $M$, killed by $\lambda^5$, is a finite locally constant
coefficient on $\operatorname{Spec}\mathcal O_K[1/(3\cdot269)]$
trivialized by its finite etale cover from $E_2$.
This is precisely the finite-cover/action dictionary of
SGA~1, Expos\'e~V, \S\S4 and 7, reused here as a categorical
interpretation.  Equations (KF.4) and (KF.12) describe the
finite specialization obtained by imposing the twisted
cyclotomic relation $4^e\gamma=1$.
Under the already cited even K-theory and etale-cohomology
comparison in AC-K10--AC-K12, this quotient is the wild
localization kernel in degree $2e$, rather than the entire
even K-group.  AC-K18 bounds its integral length and AC-K20
reduces all its specializations to one finite module with
action.  These statements supply no polarization or analytic
zero operator.  No passage in Grothendieck's handwritten
archive, or additional scan-level locator, is claimed.
